\documentclass{report}
\usepackage{amsmath,amssymb}
\setlength{\parindent}{0mm}
\setlength{\parskip}{1em}
\begin{document}
\begin{center}
\rule{6in}{1pt} \
{\large James Herring \\
{\bf Optimization Schemes for Phase Recovery in Bispectral Imaging}}

Department of Mathematics and Computer Science \\ Emory University \\ 400 Dowman Dr  \\ W401 \\ Atlanta \\ GA \\ 30322
\\
{\tt jlherri@emory.edu}\\
James Nagy\end{center}

Bispectral imaging is one class of methods used for the recovery of
images in ground-based astronomical imaging. The problem is a familiar
one: retrieve a clear, high-resolution image of an object in space given
a sequence of telescopic images of an object and a star (acting as the
point spread function or blurring operator) through an
atmosphere-telescope system. Bispectral methods approach this problem by
separating the image recovery into two subproblems: recovery of the
object�s intensity (also known as its Fourier modulus or power spectrum)
and recovery of the object�s Fourier phase. To recover the object�s
phase, we exploit the relationship between the object�s phase and the
bispectrum phase of the collected image frames, where the bispectrum is a
statistical quantity which is calculable given a sufficient number of
frames. Mathematically, this is formulated as a large-scale, non-linear
inverse problem. In this talk, we describe two previously proposed
alternative formulations of the objective function (Haniff and
Glindemann, Dainty) and discuss efficient approaches to solving these
optimization problems.

1.) A. Glindemann and J.C. Dainty, �Object fitting to the bispectral
phase by using least squares,� JOSA A 10.5 (1993): 1056-1063.

2.) C.A. Haniff, �Least-squares Fourier phase estimation from the modulo
2-pi bispectrum phase,� JOSA A 8.1 (1991): 134-140.


\end{document}
